\section{一维辅助场形式与非光锥威尔逊线算符的重正化}
\label{sec:EFT}
\subsection{一维辅助场形式}
一维辅助场形式通过在 QCD 拉格朗日量中加入如下额外项来定义 \cite{Dorn:1986dt}
\begin{eqnarray}
\label{eq:extendLBare}
\mathcal{L}_{h_v} =\bar{h}_{v, 0} (iv\cdot D) h_{v, 0},
\end{eqnarray}
其中下标“$0$”表示裸量，$h_{v, 0}$ 是 SU($3$) 规范群基础表示或伴随表示下的辅助``重''Grassmann（或复）标量场，$D_0^{\mu} = \partial^{\mu} - ig_{0}A_0^{\mu, a}T^a$ 是协变导数，$T^a$ 是相应表示下的 SU($3$) 生成元。注意，当 $v=(1,0,0,0)$ 且 $h_{v, 0}$ 处于基础表示时，eq.~(\ref{eq:extendLBare}) 的有效拉格朗日量本质上与静止系中重夸克有效理论（HQET）的拉格朗日量一致 \cite{Eichten:1989zv}，区别只在于我们取 $h_{v, 0}$ 为标量。

有效拉格朗日量 $\mathcal{L}_{h_v}$ 可通过如下重新定义进行重正化：
\begin{eqnarray}
h_{v, 0} = Z_{h_v}^{\frac{1}{2}} h_v,\, \, \, \, g_0 =Z_g g,\, \, \, \,   A_0 = Z_A^{\frac{1}{2}} A,
\end{eqnarray}
其中 $Z_{h_v}$ 是``重''辅助场的波函数重正化常数，$Z_g$ 和 $Z_A$ 是通常的强耦合与胶子场的 QCD 重正化常数。``重''辅助场会出现一个``质量修正''$i\delta m$，用重正化场和参数表达的有效拉格朗日量 $\mathcal{L}_{h_v}$ 为\footnote{在闵氏时空中，类空 $v$ 时``质量修正''$i\delta m$ 为虚数，类时 $v$ 时为实数。在欧氏时空中，$i\delta m$ 总是虚数。}
\begin{eqnarray}
\label{eq:extendLRen}
\mathcal{L}_{h_v} =Z_{h_v}\bar{h}_v (iv\cdot \partial - i\delta m) h_v + gZ_gZ_A^{\frac{1}{2}}Z_{h_v} \bar{h}_v v\cdot A^{a}T^a h_v,
\end{eqnarray}
其中对场与生成元 $T^a$ 隐含取适当的表示。

``质量修正''$i\delta m$ 是线性发散的。该线性发散在维数正则化中消失，但在格点正则化和梯度流形式中依然存在。此外，``质量修正''$i\delta m$ 可以包含 $\mathcal{O}(\Lambda_{\rm QCD})$ 量级的贡献，这来自 HQET 中已被认识到的 renormalon 歧义 \cite{Bigi:1994em, Beneke:1994sw}。

基于一维辅助场形式，``重''辅助场的连通两点函数在相差一个常数归一化因子的情况下可以通过 \cite{Dorn:1986dt,Braun:2020ymy,Ji:2020ect} 与威尔逊线相联系：
\begin{eqnarray}
\langle h_{v,0}(x)\bar{h}_{v,0}(0) \rangle_{h_v} = W\left(\frac{v\cdot x}{v^2}, 0\right)\theta\left(\frac{v\cdot x}{v^2}\right)\delta^{(d-1)}(x_\perp),
\end{eqnarray}
其中 $\langle ... \rangle_{h_v}$ 表示把``重''辅助场 $h_v$ 积出，而 $x_{\perp}^{\mu} = x^{\mu} - v^{\mu}(v\cdot x)/v^2$。
这样一来，威尔逊线算符便可以用局域流算符的乘积来替代。对于 eq.~(\ref{eq:defqoperator}) 与 eq.~(\ref{eq:defgoperator}) 定义的威尔逊线算符，我们有
\begin{eqnarray}
\mathcal{O}^{\rm B}_{\Gamma}(zv)  &=& \int d^d x\, \delta\left(\frac{v\cdot x}{v^2} -z\right)\langle \bar{\psi}_{0}(x)h_{v,0}(x) \Gamma \bar{h}_{v,0}(0)\psi_0(0)\rangle_{h_v},\\
\mathcal{O}^{\mu\nu\alpha\beta, \rm B}(zv)  &=& \int d^d x\, \delta\left(\frac{v\cdot x}{v^2} -z\right)\langle g^2F_{0}^{\mu\nu}(x)h_{v,0}(x)  \bar{h}_{v,0}(0)F_{0}^{\alpha\beta}(0)\rangle_{h_v},
\end{eqnarray}
其中上标``B''表示裸复合算符；$\mathcal{O}^{\rm B}_{\Gamma}(zv)$、$\mathcal{O}^{\mu\nu\alpha\beta, \rm B}(zv)$ 中的裸``重''辅助场 $h_{v,0}$ 分别处于基础表示和伴随表示。


\subsection{非光锥威尔逊线算符的重正化}
eq.~(\ref{eq:defgoperator}) 定义的算符不是可乘重正化的，因为在外部四矢量 $v^{\mu}$ 存在时，平行于与垂直于 $v^{\mu}$ 的分量以不同方式重正化 \cite{Braun:2020ymy}。我们引入如下投影算符
\begin{eqnarray}
g^{\mu\nu}_{\|} = \frac{v_{\mu}v_{\nu}}{v^2},\, \, \, \, g^{\mu\nu}_{\perp} = g^{\mu\nu} - \frac{v_{\mu}v_{\nu}}{v^2},
\end{eqnarray}
并通过定义
\begin{eqnarray}
F^{\mu\nu}_{\|\perp} &=& g^{\mu\alpha}_{\|}g^{\nu\beta}_{\perp}F^{\alpha\beta},\\
F^{\mu\nu}_{\perp\perp} &=& g^{\mu\alpha}_{\perp}g^{\nu\beta}_{\perp}F^{\alpha\beta}.
\end{eqnarray}
来投影出场强张量 $F^{\mu\nu}$ 的纵向与横向分量。
注意，我们忽略与规范非不变算符的混合，因为它们对规范不变可观测量没有贡献。不过，只要与我们的计算相关，我们就会验证文献 \cite{Dorn:1981wa,Zhang:2018diq,Wang:2019tgg} 中发现的混合模式。
显然，当 $v=(1,0,0,0)$ 且 $\mu\neq\nu$ 时，$F^{\mu\nu}_{\|\perp}, F^{\mu\nu}_{\perp\perp}$ 分别正比于色电场 \textbf{E} 和色磁场 \textbf{B}。我们引入如下胶子型威尔逊线算符
\begin{eqnarray}
\label{eq:defpperp}
\mathcal{O}_{\|\perp}^{\mu\nu\alpha\beta} (zv)&=& g^2 F_{\|\perp}^{\mu\nu}(zv)W(zv,0)F_{\|\perp}^{\alpha\beta}(0),\\
\label{eq:defperpperp}
\mathcal{O}_{\perp\perp}^{\mu\nu\alpha\beta} (zv)&=& g^2 F_{\perp\perp}^{\mu\nu}(zv)W(zv,0)F_{\perp\perp}^{\alpha\beta}(0).
\end{eqnarray}
如文献 \cite{Green:2017xeu, Ji:2017oey, Zhang:2018diq} 所证明、并在文献 \cite{Ishikawa:2017faj, Li:2018tpe} 中以图方法所示，eq.~(\ref{eq:defqoperator})、eq.~(\ref{eq:defpperp}) 与 eq.~(\ref{eq:defperpperp}) 定义的非光锥威尔逊线算符在辅助场形式下经历可乘重正化。也就是说，eq.~(\ref{eq:defqoperator})、eq.~(\ref{eq:defpperp}) 与 eq.~(\ref{eq:defperpperp}) 中重正化后的非光锥威尔逊线算符通过下列公式与裸算符相联系：
\begin{eqnarray}
\mathcal{O}^{\rm B}_{\Gamma}(zv, \Lambda) &=& Z_{q}(\Lambda,\mu) e^{\delta m(\Lambda)z} \mathcal{O}^{\rm R}_{\Gamma}(zv, \mu),\\
\mathcal{O}^{\mu\nu\alpha\beta,\, \rm B}_{\|\perp}(zv, \Lambda) &=& Z_{\|\perp}(\Lambda,\mu) e^{\delta m(\Lambda)z} \mathcal{O}_{\|\perp}^{\mu\nu\alpha\beta,\, \rm R}(zv, \mu),\\
\mathcal{O}^{\mu\nu\alpha\beta,\, \rm B}_{\perp\perp}(zv, \Lambda) &=& Z_{\perp\perp}(\Lambda,\mu) e^{\delta m(\Lambda)z} \mathcal{O}_{\|\perp}^{\mu\nu\alpha\beta,\, \rm R}(zv, \mu),
\end{eqnarray}
其中上标``R''表示重正化复合算符，$\mu$ 是重正化标度，$\Lambda$ 代表紫外截断（UV 调节子），$\delta m(\Lambda)$ 编码威尔逊线的``质量修正''及其线性发散，而 $Z_{q}(\Lambda,\mu)$、$Z_{\|\perp}(\Lambda,\mu)$ 与 $Z_{\perp\perp}(\Lambda,\mu)$ 编码来自波函数和顶点重正化的所有对数发散。


在一维辅助场形式中，减除线性发散之后，eq.~(\ref{eq:defqoperator})、eq.~(\ref{eq:defpperp}) 与 eq.~(\ref{eq:defperpperp}) 定义的非光锥 Wilson 线算符的重正化被简化为两个局域``重到轻''或``重到胶子''流算符的重正化\footnote{基于 Mandelstam 形式的另一方法中非光锥威尔逊线算符的重正化见文献 \cite{Stefanis:1983ke,Stefanis:2012if}。}
\begin{eqnarray}
\label{eq:currentoperator}
J_q(x) &=& \bar{\psi}(x) h_v(x),\\
J^{\mu\nu}_{\|\perp}(x) &=& gF_{\|\perp}^{\mu\nu}(x)h_v(x),\\
J^{\mu\nu}_{\perp\perp}(x) &=& gF_{\perp\perp}^{\mu\nu}(x)h_v(x).
\end{eqnarray}
上述规范不变局域流算符的重正化公式为
\begin{eqnarray}
J_q^{\rm B} (x, \Lambda) &=& Z_{J_q} (\Lambda,\mu)J_q^{\rm R}(x, \mu),\\
J^{\mu\nu,\, \rm B}_{\|\perp}(x, \Lambda) &=& Z_{J,\, \|\perp}(\Lambda,\mu)J^{\mu\nu,\, \rm R}_{\|\perp}(x, \mu),\\
J^{\mu\nu,\, \rm B}_{\perp\perp}(x, \Lambda) &=& Z_{J,\, \perp\perp}(\Lambda,\mu)J^{\mu\nu,\, \rm R}_{\perp\perp}(x, \mu),
\end{eqnarray}
由此得到等价关系
\begin{eqnarray}
Z_{q} (\Lambda,\mu)&=& Z_{J_q}^2(\Lambda,\mu),\\
Z_{\|\perp} (\Lambda,\mu)&=& Z^2_{J,\, \|\perp}(\Lambda,\mu),\\
Z_{\perp\perp} (\Lambda,\mu)&=& Z^2_{J,\, \perp\perp}(\Lambda,\mu).
\end{eqnarray}
在下文中，我们可以省略重正化常数的宗量 $\Lambda$ 与 $\mu$。

局域流算符的重正化常数还可以进一步分解为
\begin{eqnarray}
Z_{J_q} &=& Z_{\psi}^{\frac{1}{2}} Z_{h_v}^{\frac{1}{2}}Z^V_{J},\\
Z_{J,\, \|\perp} &=& Z_{A}^{\frac{1}{2}} Z_{h_v}^{\frac{1}{2}}Z^V_{\|\perp},\\
Z_{J,\, \perp\perp} &= & Z_{A}^{\frac{1}{2}} Z_{h_v}^{\frac{1}{2}}Z^V_{\perp\perp},
\end{eqnarray}
其中 $Z^V_{J}$、$Z^V_{|\perp}$ 与 $Z^V_{\perp\perp}$ 表示来自顶点修正的相应重正化因子。这一记法还隐含``重''辅助场波函数重正化常数对色表示的依赖。

以色磁场插入威尔逊圈定义的自旋依赖势所涉及的威尔逊线算符也是本研究感兴趣的对象 \cite{Brambilla:2000gk,Pineda:2000sz}。在辅助场形式中，每一次向威尔逊圈的色磁场插入都可以与如下流算符相联系
\begin{eqnarray}
\label{eq:currentoperatorspin}
J^{\mu\nu}_{F}(x) &=&\bar{h}_v(x)gF^{\mu\nu}_{\perp\perp}(x)h_v(x),
\end{eqnarray}
其中 $v=(1, 0,0,0)$，下标 $F$ 表示``重''辅助场处于 SU($3$) 群的基础表示。该算符的重正化公式为
\begin{eqnarray}
J^{\mu\nu,\, \rm B}_{F}(x)&=&Z_{F}J^{\mu\nu,\, \rm R}_{F}(x) = Z_{A}^{1/2} Z_{h_v}Z^V_{F}J^{\mu\nu,\, \rm R}_{F}(x),
\end{eqnarray}
其中 $Z^V_{F}$ 表示来自流 $\bar{h}_v(x)gF^{\mu\nu}_{\perp\perp}(x)h_v(x)$ 顶点修正的重正化因子。
